\documentclass[12pt]{article}
\usepackage[utf8]{inputenc}
\usepackage[T1]{fontenc}
\usepackage[margin=1in]{geometry}
\usepackage{setspace}
\usepackage{xcolor}
\usepackage[colorlinks=true, urlcolor=blue!60!black, linkcolor=blue!60!black]{hyperref}
\setlength{\parindent}{0em}
\setlength{\parskip}{0.45em}
\renewcommand{\baselinestretch}{1.0}
\begin{document}
\begin{center}
{\Large \textit{TransUnion LLC v. Ramirez}, 594 U.S. 413 (2021)\par}
\vspace{0.3em}
{\normalsize Case Brief --- One-Page Edition\par}
\end{center}

\textbf{Memory Jogger.} No concrete harm, no standing: an unsent defamatory letter in a desk drawer harms no one (\href{https://supreme.justia.com/cases/federal/us/594/20-297/}{Justia}).

\textbf{Facts.} TransUnion falsely flagged 8,185 consumers as potential matches to a terrorist list; only 1,853 had reports sent to businesses. A dealer refused to sell Sergio Ramirez a car because of the alert, and TransUnion's follow-up mailings omitted required information.

\textbf{Procedural History.} A jury awarded the class over \$60 million; the Ninth Circuit found standing for all.

\textbf{Votes.} Reversed 5--4: Kavanaugh for the Court; Thomas and Kagan dissented.

\textbf{Holding.} Each class member seeking damages must show concrete harm with a close common-law analogue; only those whose reports were disseminated had standing.

\textbf{Reasoning.} Congress cannot enact an injury into existence. Disseminated alerts resemble defamation; internal files do not; unmaterialized risks cannot support damages, though they may support an injunction. Thomas would let any violation of a private right suffice; Kagan saw judicial aggrandizement.

\textbf{Examples.} \emph{Same-Side:} retained but undisclosed data. \emph{Opposite-Side:} viewing histories sent to advertisers. \emph{Fence-Sitter:} data sent to a print vendor.

\textbf{Critique.} The analogue test is indeterminate and pushes class actions into state courts.

\textbf{Key Quotation.} ``Congress may enact legal prohibitions and obligations. And Congress may create causes of action for plaintiffs to sue defendants who violate those legal prohibitions or obligations. But under Article III, an injury in law is not an injury in fact. Only those plaintiffs who have been concretely harmed by a defendant's statutory violation may sue that private defendant over that violation in federal court.'' --- Justice Kavanaugh, 594 U.S. at 427.

\end{document}
